\documentclass[11pt,a4paper,fontset=fandol]{ctexart}

\usepackage[a4paper,top=2.4cm,bottom=2.6cm,left=2.35cm,right=2.35cm,headheight=18pt,headsep=0.55cm,footskip=26pt]{geometry}
\usepackage{amsmath,amssymb}
\usepackage{fancyhdr}
\usepackage{lastpage}
\usepackage{enumitem}
\usepackage{titlesec}
\usepackage{xcolor}
\usepackage{hyperref}
\usepackage{bookmark}
\usepackage[most]{tcolorbox}
\usepackage{setspace}
\usepackage{array}

\hypersetup{
  colorlinks=true,
  linkcolor=blue!50!black,
  urlcolor=blue!50!black
}

\setstretch{1.16}
\setlength{\parindent}{2em}
\setlength{\parskip}{0.32em}
\setlist[enumerate]{itemsep=0.35em, topsep=0.35em, leftmargin=2.3em}
\setlist[itemize]{itemsep=0.25em, topsep=0.25em, leftmargin=2em}

\definecolor{mainblue}{RGB}{36,83,140}
\definecolor{lightblue}{RGB}{241,246,252}
\definecolor{lightgray}{RGB}{246,246,246}
\definecolor{rulegray}{RGB}{110,118,128}
\definecolor{softgreen}{RGB}{236,247,240}
\definecolor{softyellow}{RGB}{254,249,229}

\titleformat{\section}
  {\Large\bfseries\color{mainblue}}
  {\thesection.}
  {0.45em}
  {}
  [\vspace{0.25em}\color{rulegray}\titlerule]
\titlespacing*{\section}{0pt}{1.2em}{0.8em}
\titleformat{\subsection}{\large\bfseries\color{mainblue}}{\thesubsection}{0.5em}{}

\pagestyle{fancy}
\fancyhf{}
\fancyhead[L]{\small 数论方法中的组合构造专题目录页}
\fancyhead[C]{\small 组合专题资料索引}
\fancyhead[R]{\small 刘文博}
\fancyfoot[C]{\small 第 \thepage 页 / 共 \pageref{LastPage} 页}
\renewcommand{\headrulewidth}{0.55pt}
\renewcommand{\footrulewidth}{0.35pt}

\newtcolorbox{goalbox}[1]{
  breakable,
  colback=lightblue,
  colframe=mainblue,
  boxrule=0.7pt,
  arc=1mm,
  title=\textbf{#1},
  fonttitle=\bfseries
}

\newtcolorbox{infobox}[1]{
  breakable,
  colback=softgreen,
  colframe=green!40!black,
  boxrule=0.65pt,
  arc=1mm,
  title=\textbf{#1},
  fonttitle=\bfseries
}

\newtcolorbox{warnbox}[1]{
  breakable,
  colback=softyellow,
  colframe=orange!60!black,
  boxrule=0.65pt,
  arc=1mm,
  title=\textbf{#1},
  fonttitle=\bfseries
}

\begin{document}

\thispagestyle{empty}
\begin{titlepage}
\centering
\vspace*{0.9cm}

\begin{tcolorbox}[
  enhanced,
  width=0.92\textwidth,
  colback=white,
  colframe=mainblue,
  boxrule=1pt,
  arc=0mm,
  left=6mm,
  right=6mm,
  top=8mm,
  bottom=8mm
]
\centering

{\fontsize{24}{30}\selectfont\bfseries 数论方法中的组合构造专题目录页\par}
\vspace{0.65cm}
{\fontsize{17}{22}\selectfont 讲义、题单与周测的使用说明\par}

\vspace{1cm}
\begin{tcolorbox}[colback=lightblue,colframe=mainblue,width=0.84\textwidth,boxrule=1pt]
\centering
{\large 适合对象：已经掌握基础计数、准备把“奇偶性和余数分类”真正用进组合题的学生}\\[0.35em]
{\normalsize 本目录页帮助把“数论工具进入组合证明”的资料串成一个完整训练包}
\end{tcolorbox}

\vspace{1.0cm}
\renewcommand{\arraystretch}{1.42}
\begin{tabular}{>{\bfseries}p{3.5cm}p{8cm}}
专题名称： & 数论方法中的组合构造 \\
资料组成： & 课堂讲义、提高训练题单、周测 A 卷、周测 B 卷 \\
使用方式： & 先学讲义，再做题单，最后用 A/B 周测检查掌握情况 \\
编译方式： & XeLaTeX \\
\end{tabular}

\vspace{0.9cm}
{\large \today\par}
\end{tcolorbox}
\end{titlepage}

\tableofcontents
\newpage

\section{专题包包含哪些资料}

\begin{goalbox}{四份核心资料}
\begin{enumerate}
  \item 课堂讲义：帮助理解奇偶性、模运算、余数分类、不变量与基础构造；
  \item 提高训练题单：用于课后巩固和变式提升，重点训练“会分类、会证明、会做操作题”；
  \item 周测 A 卷：检查基础按模分类和不变量是否稳定；
  \item 周测 B 卷：检查综合构造、存在性证明和较难分类题是否真正掌握。
\end{enumerate}
\end{goalbox}

\begin{infobox}{对应文件名}
\begin{itemize}
  \item 讲义：\texttt{number\_theory\_constructions\_handout.tex} / \texttt{number\_theory\_constructions\_handout.pdf}
  \item 提高训练题单：\texttt{number\_theory\_constructions\_advanced\_problem\_set.tex} / \texttt{number\_theory\_constructions\_advanced\_problem\_set.pdf}
  \item 周测 A 卷：\texttt{number\_theory\_constructions\_weekly\_test\_A.tex} / \texttt{number\_theory\_constructions\_weekly\_test\_A.pdf}
  \item 周测 B 卷：\texttt{number\_theory\_constructions\_weekly\_test\_B.tex} / \texttt{number\_theory\_constructions\_weekly\_test\_B.pdf}
\end{itemize}
\end{infobox}

\section{建议学习顺序}

\begin{goalbox}{推荐顺序}
\begin{enumerate}
  \item 先完整学习讲义，重点吃透“为什么只看奇偶性就能判断不可能”“为什么按余数分类能做存在性证明”；
  \item 再完成提高训练题单，训练把“看起来是选数题”的问题转化成模 $2$、模 $3$、模 $4$ 的分类；
  \item 接着做周测 A 卷，检查基础按模分类与简单不变量是否稳定；
  \item 最后做周测 B 卷，检查较难构造、操作题与综合证明是否掌握。
\end{enumerate}
\end{goalbox}

\begin{warnbox}{不建议的做法}
\begin{itemize}
  \item 一看到数论味道就盲目上同余式，而不先判断“按模几分类最自然”；
  \item 只会验证一个构造例子，却不会说明为什么所有情况都只能这样；
  \item 操作题只跟着做，不先找“每一步都不变”的量；
  \item 证明题里把“至少存在一个”误写成“任意一个都如此”。
\end{itemize}
\end{warnbox}

\section{每份资料主要解决什么问题}

\subsection{课堂讲义}
讲义适合第一次系统学习这个专题，主要解决以下问题：
\begin{itemize}
  \item 什么情况下先看奇偶性就够了；
  \item 什么情况下应按模 $3$、模 $4$、模 $5$ 分类；
  \item 操作型问题里不变量该怎么找；
  \item 棋盘、格点、选数分组这些题里，数论工具怎样转化成组合限制。
\end{itemize}

\subsection{提高训练题单}
题单适合第二阶段训练，主要用来把方法做熟。它更强调：
\begin{itemize}
  \item 能否快速判断该按什么模分类；
  \item 能否完整列出所有可能余数结构；
  \item 能否证明“最多”“至少”“不可能”这类结论；
  \item 能否从操作中识别出总和、奇偶性、余数等不变量。
\end{itemize}

\subsection{周测 A 卷}
A 卷偏基础，适合检查：
\begin{itemize}
  \item 奇偶性与余数分类是否熟练；
  \item 简单构造题是否会做；
  \item 抽屉 + 按模分类是否能结合使用；
  \item 基础不变量题是否会判断可能与不可能。
\end{itemize}

\subsection{周测 B 卷}
B 卷偏综合，适合检查：
\begin{itemize}
  \item 分类是否完整；
  \item 证明是否覆盖所有情况；
  \item 会不会把“模运算”和“组合结构”结合起来；
  \item 能否自己构造一个满足条件的分组或反例。
\end{itemize}

\section{一个推荐的 4 次训练安排}

\begin{goalbox}{4 次完成一个小专题}
\begin{enumerate}
  \item 第 1 次：学习讲义前半部分，重点理解奇偶性、模运算和基础例题；
  \item 第 2 次：学习讲义后半部分，完成题单基础题和中档题；
  \item 第 3 次：完成题单综合题，并把错题按“模数选错/分类漏了/不变量没抓到/构造不完整”分类订正；
  \item 第 4 次：完成周测 A 卷或 B 卷，作为阶段性检查。
\end{enumerate}
\end{goalbox}

\begin{infobox}{如果孩子基础较好，可以这样压缩}
\begin{itemize}
  \item 第 1 天：讲义 + 题单基础题；
  \item 第 2 天：题单变式题 + 错题订正；
  \item 第 3 天：周测 A 卷；
  \item 第 4 天：周测 B 卷。
\end{itemize}
\end{infobox}

\section{专题学习后的自检清单}

\begin{goalbox}{如果下面 8 条都能做到，说明这个专题基本过关}
\begin{enumerate}
  \item 我会判断一道题更适合看奇偶性还是看模 $3$、模 $4$、模 $5$；
  \item 我会把一个集合按余数类完整分类；
  \item 我会列出若干个数和为某个模数的所有余数结构；
  \item 我会证明“至少存在一个”类型的结论；
  \item 我会证明“不可能”类型的结论；
  \item 我会找到一些简单操作题里的不变量；
  \item 我会构造一个满足和为 $3$ 的倍数或 $4$ 的倍数的分组；
  \item 我会在解答后检查分类是否完整、论证是否遗漏情况。
\end{enumerate}
\end{goalbox}

\section{给家长的使用建议}

\begin{warnbox}{更有效的带练方式}
\begin{itemize}
  \item 不要一开始就提示“用模运算”，先让孩子自己说“这题该按什么分”；
  \item 如果卡住，可以只提示“先把数按余数分组看看”；
  \item 操作题订正时，重点看孩子有没有主动寻找不变量；
  \item 做完后请孩子口头复述：这题到底是用奇偶性、抽屉、余数分类，还是不变量解决的。
\end{itemize}
\end{warnbox}

\begin{infobox}{和后续专题的衔接}
这个专题学完以后，最自然的下一步是继续学习\textbf{抽屉原理与极值思想}以及\textbf{染色方法}。因为很多存在性证明、不可能性证明和棋盘问题，都会和按模分类、不变量思想自然结合起来。
\end{infobox}

\end{document}
